\section{Architektura systemu}

\subsection{Przegląd architektury}

System symulacji pożarów lasu zbudowany jest w architekturze mikroserwisów, składającej się z czterech głównych modułów oraz infrastruktury wspierającej:

\begin{itemize}
    \item \textbf{FFBackend} (fire-backend) -- serwis Java/Spring Boot odpowiedzialny za logikę biznesową i koordynację
    \item \textbf{FFSim} (fire-simulation) -- silnik symulacji napisany w Pythonie
    \item \textbf{FFSup} (fire-support) -- serwis wsparcia decyzyjnego wykorzystujący MCTS
    \item \textbf{FFVis} (fire-visualization) -- aplikacja frontendowa do wizualizacji
    \item \textbf{FFConf} (fire-configurations) -- serwis zarządzania konfiguracjami lasu
\end{itemize}


\subsection{Technologie i narzędzia}

\begin{itemize}
    \item \textbf{Backend:} Java 17, Spring Boot, Maven
    \item \textbf{Symulacja:} Python 3.x, asyncio
    \item \textbf{Frontend:} TypeScript, React, MapLibre GL
    \item \textbf{Komunikacja:} RabbitMQ (topic exchange), HTTP/REST, Server-Sent Events (SSE)
    \item \textbf{Baza danych:} MongoDB
    \item \textbf{Orchestracja:} Docker, Docker Compose
    \item \textbf{Logowanie:} Grafana, Loki, Promtail (w trakcie implementacji)
\end{itemize}

\section{Komponenty systemu}

\subsection{FFBackend (fire-backend)}

\subsubsection{Odpowiedzialności}

\begin{itemize}
    \item Zarządzanie stanem symulacji (sektory, brygady, patrole)
    \item Agregacja danych z symulacji
    \item Udostępnianie REST API dla frontendu
    \item Strumieniowanie stanu symulacji przez SSE
    \item Przekazywanie komend z frontendu do symulacji
    \item Publikowanie zagregowanych danych do serwisu wsparcia
\end{itemize}


\subsubsection{Zależności}

\begin{itemize}
    \item Wymaga: RabbitMQ, MongoDB (opcjonalnie)
    \item Komunikuje się z: FFSim (HTTP + RabbitMQ), FFSup (RabbitMQ), FFVis (HTTP/SSE)
    \item Port: 8181 (domyślnie)
\end{itemize}

\subsection{FFSim (fire-simulation)}

\subsubsection{Odpowiedzialności}

\begin{itemize}
    \item Wykonanie symulacji fizycznej pożaru lasu
    \item Modelowanie rozprzestrzeniania się ognia
    \item Zarządzanie agentami (brygady pożarnicze, patrole leśne)
    \item Aktualizacja stanu sektorów (poziom ognia, spalenia, gaszenia)
    \item Publikowanie danych telemetrycznych (czujniki, agenci, sektory)
    \item Odbieranie komend sterujących od backendu
\end{itemize}


\subsubsection{Zależności}

\begin{itemize}
    \item Wymaga: RabbitMQ
    \item Komunikuje się z: FFBackend (HTTP + RabbitMQ)
    \item Port: 5000 (domyślnie)
\end{itemize}

\subsubsection{Ważna uwaga}

Symulator jest osobnym bytem w systemie. Wszystkie decyzje podejmowane są przez system (FFBackend, FFSup), a symulator jedynie wykonuje komendy i symuluje fizykę pożaru.

\subsection{FFSup (fire-support)}

\subsubsection{Odpowiedzialności}

\begin{itemize}
    \item Generowanie rekomendacji decyzyjnych przy użyciu MCTS
    \item Analiza stanu symulacji
    \item Publikowanie rekomendacji do backendu/frontendu
    \item Zarządzanie pulą workerów MCTS
    \item Konwersja stanu symulacji do formatu wymaganego przez algorytmy
\end{itemize}

\subsubsection{Zależności}

\begin{itemize}
    \item Wymaga: RabbitMQ
    \item Komunikuje się z: FFBackend (RabbitMQ)
    \item Port: brak (tylko RabbitMQ)
\end{itemize}

\subsection{FFVis (fire-visualization)}

\subsubsection{Odpowiedzialności}

\begin{itemize}
    \item Wizualizacja mapy lasu i sektorów
    \item Wyświetlanie pozycji brygad i patroli
    \item Prezentacja rekomendacji decyzyjnych
    \item Interakcja użytkownika (wysyłanie komend)
    \item Wyświetlanie danych z czujników
\end{itemize}

\subsubsection{Zależności}

\begin{itemize}
    \item Wymaga: FFBackend (HTTP/SSE)
    \item Port: 8080 (domyślnie)
\end{itemize}

\subsection{FFConf (fire-configurations)}

\subsubsection{Odpowiedzialności}

\begin{itemize}
    \item Zarządzanie konfiguracjami lasu w MongoDB
    \item Udostępnianie API do CRUD konfiguracji
    \item Walidacja konfiguracji JSON
    \item Przechowywanie hierarchicznej struktury konfiguracji
\end{itemize}


\subsubsection{Zależności}

\begin{itemize}
    \item Wymaga: MongoDB
    \item Port: 31415 (domyślnie)
\end{itemize}

\section{Komunikacja między modułami}

\subsection{Protokoły komunikacji}

\subsubsection{HTTP/REST}

Używany do:
\begin{itemize}
    \item Komunikacji FFVis $\leftrightarrow$ FFBackend
    \item Komunikacji FFBackend $\leftrightarrow$ FFSim (komendy, konfiguracja)
    \item Komunikacji FFVis $\leftrightarrow$ FFConf (pobieranie konfiguracji)
\end{itemize}

\subsubsection{RabbitMQ (Topic Exchange)}

Używany do:
\begin{itemize}
    \item Publikacji danych telemetrycznych z FFSim do FFBackend
    \item Publikacji zagregowanych danych z FFBackend do FFSup
    \item Publikacji rekomendacji z FFSup do FFBackend
    \item Kolejkowania komend gdy FFSim jest niedostępny
\end{itemize}

\subsubsection{Server-Sent Events (SSE)}

Używany do:
\begin{itemize}
    \item Strumieniowania stanu symulacji z FFBackend do FFVis
    \item Jednokierunkowa komunikacja w czasie rzeczywistym
\end{itemize}

\subsection{Exchange RabbitMQ}

Wszystkie kolejki używają exchange typu \texttt{topic} o nazwie \texttt{fire-simulation-exchange} (konfigurowalne przez zmienną środowiskową \texttt{RABBITMQ\_EXCHANGE}).

\section{Uruchomienie systemu}

\subsection{Wymagania}

\begin{itemize}
    \item Docker i Docker Compose
    \item Minimum 4GB RAM
    \item Porty: 8080, 8181, 5000, 31415, 27017, 5672, 15672
\end{itemize}

\subsection{Uruchomienie przez Docker Compose}

\subsubsection{Uruchomienie wszystkich serwisów}

\begin{lstlisting}[language=bash, basicstyle=\small]
docker compose -f docker-compose.yml up --build
\end{lstlisting}

\subsubsection{Uruchomienie wybranych serwisów}

\begin{lstlisting}[language=bash, basicstyle=\small]
# Najpierw infrastruktura
docker compose up --build rabbitmq-service mongo

# Następnie serwisy aplikacyjne
docker compose up --build fire-backend-service
docker compose up --build fire-simulation-service
docker compose up --build fire-configuration-service
docker compose up --build fire-visualization-service
\end{lstlisting}

\subsection{Kolejność uruchamiania}

Zalecana kolejność uruchamiania serwisów:

\begin{enumerate}
    \item \textbf{MongoDB} -- baza danych dla konfiguracji
    \item \textbf{RabbitMQ} -- broker komunikatów
    \item \textbf{FFConf} -- serwis konfiguracji (wymaga MongoDB)
    \item \textbf{FFBackend} -- backend (wymaga RabbitMQ, opcjonalnie MongoDB)
    \item \textbf{FFSim} -- symulator (wymaga RabbitMQ)
    \item \textbf{FFSup} -- serwis wsparcia (wymaga RabbitMQ)
    \item \textbf{FFVis} -- frontend (wymaga FFBackend)
\end{enumerate}

\subsection{Sprawdzanie statusu}

\begin{itemize}
    \item \textbf{FFBackend:} \texttt{http://localhost:8181/actuator/health}
    \item \textbf{FFConf:} \texttt{http://localhost:31415/health}
    \item \textbf{RabbitMQ Management:} \texttt{http://localhost:15672}
    \item \textbf{FFVis:} \texttt{http://localhost:8080}
\end{itemize}

\section{Architektura agentowa}

\subsection{Model agentów}

System wykorzystuje architekturę agentową, gdzie agenci reprezentują autonomiczne jednostki działające w symulacji.

\subsubsection{Definicja agenta}

Agent w systemie to:
\begin{itemize}
    \item \textbf{Instancja runtime:} obiekt w pamięci reprezentujący jednostkę (brygada pożarnicza, patrol leśny)
    \item \textbf{Cykl życia:} inicjalizacja $\rightarrow$ decyzja $\rightarrow$ akcja $\rightarrow$ raport
    \item \textbf{Miejsce w architekturze:} agenci "żyją" w FFSim, ale są zarządzani przez FFBackend
\end{itemize}

\subsubsection{Typy agentów}

\begin{itemize}
    \item \textbf{FireBrigade} -- brygada pożarnicza
    \begin{itemize}
        \item Stany: AVAILABLE, TRAVELLING, EXTINGUISHING
        \item Zadania: gaszenie pożarów w sektorach
        \item Decyzje: wybór sektora do gaszenia (może być LLM-driven)
    \end{itemize}
    
    \item \textbf{ForesterPatrol} -- patrol leśny
    \begin{itemize}
        \item Stany: AVAILABLE, TRAVELLING, PATROLLING
        \item Zadania: patrolowanie sektorów, wykrywanie zagrożeń
        \item Decyzje: wybór sektora do patrolowania (heurystyka, może być LLM-driven)
    \end{itemize}
\end{itemize}

\subsection{Separacja logiki decyzyjnej}

\textbf{Ważne:} Agent jest oddzielony od logiki decyzyjnej. Logika decyzyjna może być:
\begin{itemize}
    \item Deterministyczna (reguły, heurystyki)
    \item MCTS (Monte Carlo Tree Search)
    \item LLM-driven (wymienna strategia decyzyjna)
\end{itemize}

LLM jest traktowany jako \textbf{jedna z implementacji} modułu decyzyjnego.

\section{Logi systemu}

\subsection{Struktura logów}

Wszystkie serwisy logują do plików w katalogu \texttt{logs/}:
\begin{itemize}
    \item \texttt{logs/fire-backend/application.log}
    \item \texttt{logs/fire-simulation/application.log}
    \item \texttt{logs/fire-configuration/application.log}
\end{itemize}

\section{Wersjonowanie i kompatybilność}

\subsection{Wersjonowanie komunikatów}

Obecnie system nie implementuje formalnego wersjonowania komunikatów. Zalecane jest dodanie wersji do:
\begin{itemize}
    \item Kontraktów RabbitMQ (pole \texttt{version} w nagłówkach)
    \item REST API (prefiks \texttt{/api/v1/})
    \item Formatów JSON konfiguracji
\end{itemize}

\section{Ograniczenia systemu}

\subsection{System alertów}

System \textbf{nie posiada} mechanizmu alertów w czasie rzeczywistym. W przyszłości można dodać koncepcyjny mechanizm alertów jako osobny moduł.

\subsection{Charakter symulacyjny}

Wszystkie dane w systemie są generowane przez symulator. System nie korzysta z:
\begin{itemize}
    \item Realnych czujników środowiskowych
    \item Realnych danych meteorologicznych (oprócz symulowanego wiatru)
    \item Realnych jednostek pożarniczych
\end{itemize}

\section{Podsumowanie}

System symulacji pożarów lasu jest zbudowany w architekturze mikroserwisów z wyraźnym podziałem odpowiedzialności. 

